\documentclass[10pt]{article}
\usepackage[T1]{fontenc}
\usepackage[utf8]{inputenc}
\usepackage[french]{babel}
\usepackage{amsmath,amssymb,mathrsfs,mathpazo}
\usepackage{geometry}
\geometry{a4paper,margin=2cm}
\usepackage{graphicx,array,tabularx,tabularray,colortbl,multirow,multicol,booktabs,enumitem,pifont,ulem,fancybox,diagbox,cancel}
\usepackage{pstricks,pstricks-add,pst-plot,pst-tree,pst-node,pst-text,pst-func,pst-math,pst-bspline,pst-3dplot,pst-eucl,pst-solides3d}
\let\pscalculate\undefined
\usepackage{pst-bezier}
\usepackage{tikz,pgfplots,tkz-tab}
\usetikzlibrary{arrows,arrows.meta,patterns,calc,positioning,trees,decorations.markings,shapes.geometric,shapes.misc,shapes.arrows,intersections,decorations.pathmorphing,decorations.pathreplacing,backgrounds,shadows,babel}
\pgfplotsset{compat=1.18}
\usepgfplotslibrary{fillbetween}
\usepackage[np]{numprint}
\usepackage{listings,makecell,pas-tableur,twoopt}
\let\barre\undefined
\usepackage[pstricks]{tablvar}
\definecolor{darkgreen}{RGB}{0,100,0}
\providecommand{\red}{\color{red}}
\providecommand{\blue}{\color{blue}}
\providecommand{\green}{\color{green}}
\providecommand{\black}{\color{black}}
\providecommand{\white}{\color{white}}
\providecommand{\gray}{\color{gray}}
\providecommand{\cyan}{\color{cyan}}
\providecommand{\magenta}{\color{magenta}}
\providecommand{\yellow}{\color{yellow}}
\providecommand{\violet}{\color{violet}}
\providecommand{\orange}{\color{orange}}
\providecommand{\euro}{\texteuro}
\providecommand{\EUR}{\texteuro}
\providecommand{\ieme}{\textsuperscript{e}}
\providecommand{\degres}{\ensuremath{^\circ}}
\providecommand{\degree}{\ensuremath{^\circ}}
\providecommand{\vv}[1]{\overrightarrow{#1}}
\providecommand{\Coord}[1]{#1}
\providecommand{\decosix}{}
\providecommand{\starredbullet}{\ensuremath{\bullet}}
\pagestyle{empty}
\setlength{\parindent}{0pt}
\usepackage{adjustbox}\DeclareUnicodeCharacter{2236}{:}\DeclareUnicodeCharacter{2190}{\ensuremath{\leftarrow}}
\providecommand{\R}{\mathbb{R}}
\providecommand{\N}{\mathbb{N}}
\providecommand{\Z}{\mathbb{Z}}
\providecommand{\Q}{\mathbb{Q}}
\providecommand{\D}{\mathbb{D}}
\providecommand{\C}{\mathbb{C}}
\providecommand{\e}{{\mathrm{e}}}
\providecommand{\vect}[1]{\overrightarrow{#1}}
\providecommand{\vectt}[1]{\overrightarrow{\mathrm{#1}}}
\providecommand{\vvt}[1]{\overrightarrow{\mathrm{#1}}}
\providecommand{\np}[2][]{{#2~#1}}
\providecommand{\vv}[1]{{\overrightarrow{#1}}}
\renewcommand{\d}{{\mathrm{d}}}
\providecommand{\degres}{{^\circ}}
\providecommand{\degre}{{^\circ}}
\providecommand{\blue}{\color{blue}}
\providecommand{\red}{\color{red}}
\providecommand{\enskip}{\quad}
\providecommand{\llbracket}{\lbrack\!\lbrack}
\providecommand{\rrbracket}{\rbrack\!\rbrack}
\providecommand{\euro}{{\text{€}}}
\providecommand{\no}{{\text{n°}}}
\providecommand{\interval}[2]{{[#1;#2]}}
\providecommand{\intervFO}[2]{{[#1;#2[}}
\providecommand{\dotfill}{\ldots}

\providecommand{\ds}{\displaystyle}
\providecommand{\textding}[1]{#1}
\providecommand{\Oij}{\(\mathrm{O};\vec{\imath},\vec{\jmath}\)}
\providecommand{\Oijk}{\(\mathrm{O};\vec{\imath},\vec{\jmath},\vec{k}\)}
\providecommand{\Ouv}{\(\mathrm{O};\vec{u},\vec{v}\)}
\providecommand{\intervFF}[2]{\ensuremath{[#1;#2]}}
\renewcommand{\euro}{\ifmmode\text{\texteuro}\else\texteuro\fi}
\DeclareUnicodeCharacter{00B0}{\ifmmode^{\circ}\else\textdegree\fi}
\DeclareUnicodeCharacter{20AC}{\ifmmode\text{\texteuro}\else\texteuro\fi}
% Énoncés : Éducation nationale. Transcription et corrigés : APMEP, Denis Vergès et contributeurs. Découpage et classement : CoopMaths, Sébastien Lozano, Éric Elter et contributeurs. Adaptation de la présentation : Maths Explorer.
% CC BY-SA : https://coopmaths.fr/www/about/
% https://coopmaths.fr/alea/static/bac/2026/tex/bac_2026_05_sujet2_ameriquenord_3_cor.tex
\begin{document}
\noindent\textit{APMEP / CoopMaths. CC BY-SA. Présentation adaptée par Maths Explorer.}\par\medskip
L'espace est muni d'un repère orthonormé \Oijk.

On considère :

\begin{itemize}[label=$\bullet~$]
\item les points A(4~;~2~;~2), B$(5~;~-2~;~3)$ et C(1~;~1~;~1) ;
\item la droite $\Delta$ dont une représentation paramétrique est donnée par

\[\left\{\begin{array}{l c l}
x &=& 1 + 2t\\
y &=&1+ \phantom{2}t\\
z &=& 1 +2t\\
\end{array}\right. \text{avec}\: t \in \R\]
	
	\item le plan $\mathcal{P}$ passant par le point A et perpendiculaire à la droite $\Delta$.
	\end{itemize}
	
	\medskip
	
	\begin{enumerate}
		\item \begin{itemize} \item Pour $t=0$ on trouve: 

			$\begin{cases}  x= 1+2\times 0 =1 =x_C \\ y=1 +0 =1=y_C\\ z=1 +2\times 0 =1=z_C \end{cases}$ 

			donc le point C appartient à la droite $\Delta$. 

			\item Le point A appartient à la droite $\Delta$  si et seulement si il existe un réel $t$ tel que $\begin{cases} x_A=1+2t \\ y_A=1+t \\ z_A=1+2t \end{cases} $ 

c'est à dire si et seulement si le système  $\begin{cases} 4=1+2t \\ 2=1+t \\ 2=1+2t \end{cases}$ admet une unique solution.

Or la première équation donne $t=\dfrac{3}{2}$ et la deuxième $t=1$, ce système n'admet donc pas de solution d'où la droite $\Delta$ ne passe pas par le point A.
		\end{itemize} 
		\item
		\begin{enumerate}
			\item Le plan $\mathcal{P}$  perpendiculaire à la droite $\Delta$, un vecteur directeur de la droite $\Delta$ est donc un vecteur normal à $\mathcal{P}$. 
			
			La représentation paramétrique de $\Delta$ nous  donne $\vect{n}\begin{pmatrix} 2 \\1\\2 \end{pmatrix}$. 
			
			Une équation cartésienne du plan $\mathcal{P}$ est donc de la forme $2x+y+2x+d=0$ avec $d\in \R$. 
			
			De plus, le plan $\mathcal{P}$ passe  par le point A donc : 
			
			$2x_A+y_A+z_A+d=0 \Longleftrightarrow 2 \times 4 +2 +2 \times 2 +d=0 \Longleftrightarrow 8 +2 +4 +d=0 \Longleftrightarrow d=-14$

			Une équation cartésienne de $\mathcal{P}$ est donc bien :   $2x + y + 2z - 14 = 0$.
			\item  Un point appartient à un plan si et seulement si ses coordonnées vérifient une équation de ce plan.

			$2x_B+y_B+z_B-14 = 2 \times 5 -2 +2 \times 3 -14 =10-2+6-14 =0 $ 

			donc le point B appartient au plan $\mathcal{P}$.
			
			$2x_C+y_C+z_C-14 = 2 \times 1 +1 +2 \times 1 -14 = 2+1+2-14=-9 \neq 0 $ 
			
			donc le point C n'appartient pas au plan $\mathcal{P}$. 
			
		\end{enumerate}
		\item 
		\begin{enumerate}
			\item    Le point D est le projeté orthogonal du point C sur le plan $\mathcal{P}$ si et seulement si le point D appartient au plan $\mathcal{P}$ et si la droite (DC) est orthogonale au plan $\mathcal{P}$.

On a $\vectt{CD} \begin{pmatrix} 2\\1\\2  \end{pmatrix}$ et on sait que $\vect{n}\begin{pmatrix} 2 \\1\\2 \end{pmatrix}$ est un vecteur normal au plan $\mathcal{P}$.

Les vecteurs $\vectt{CD}$ et $\vect{n}$ sont égaux donc la droite (CD) est orthogonale au plan $\mathcal{P}$

De plus : $2x_D+y_D+z_D-14 = 2 \times 3 +2 +2 \times 3 -14 =6 +2+6-14 =0 $

donc le point D appartient au plan $\mathcal{P}$.

Donc le point D est le projeté orthogonal du point C sur le plan $\mathcal{P}$.
	\item  Par définition du point A, il appartient au plan $\mathcal{P}$. 

D'après la question 2.b., le point B appartient au plan $\mathcal{P}$

D est le projeté du point C sur le plan $\mathcal{P}$, donc il appartient au plan $\mathcal{P}$.

Les points A, B  et D  appartiennent donc tous les trois au plan $\mathcal{P}$  donc $\mathcal{P} = (ABD)$
	
or le point C n'appartient pas au plan  $\mathcal{P}$ d'après la question 2. b., donc les points A, B, C et D ne sont pas coplanaires.
			
			\item On a $\vectt{AB} \begin{pmatrix} 1\\-4\\1  \end{pmatrix}$ et $\vectt{AD} \begin{pmatrix} -1\\0\\1 \end{pmatrix}$

On est dans un repère orthonormé donc : 

$\vect{\text{AB}} \cdot \vect{\text{AD}}= 1 \times (-1) -4 \times 0 + 1 \times 1 =0$.
		\item Le produit scalaire $\vect{\text{AB}} \cdot \vect{\text{AD}}$ est nul donc ABD est un triangle rectangle en A.

			Soit $\mathcal{A}$ l'aire du triangle ABD. 

			On a $\begin{aligned}[t] \mathcal{A} &= \dfrac{\mathrm{AB} \times \mathrm{BD}}{2}\\ 
	& = \dfrac{ \sqrt{1^2+4^2+1^2}\times \sqrt{(-1)^2+0^2+1^2}}{2}\\
	& = \dfrac{\sqrt{18} \times \sqrt{2}}{2}\\
	& = \dfrac{\sqrt{2 \times 9 \times 2}}{2}\\
	&= 3 
	\end{aligned} $.

Dans la pyramide ABCD, la hauteur associé à la base ABD sera le segment [CD] car le point D est le projeté orthogonal du point C sur le plan (ABD). 

D'où $\begin{aligned}[t]
V&= \dfrac{1}{3} \times 3 \times \mathrm{CD}\\
&= \sqrt{2^2+1^2+2^2}\\
&=\sqrt{9}=3
\end{aligned}$ 

Le volume de la pyramide ABCD est 3 unités de volume. 
		\end{enumerate}
					\item H est le projeté orthogonal du point A sur la droite (BC) si et seulement H appartient à la droite (BC) et les droites (BC) et (AH) sont perpendiculaires.

\begin{itemize}
\item Soit $X\left(\dfrac{73}{29}~;~\dfrac{- 4}{29}~;~\dfrac{51}{29} \right)$.

On a $\vectt{BC} \begin{pmatrix} -4 \\3 \\ -2 \end{pmatrix}$ et $\vectt{\text{AX}} \begin{pmatrix} -\frac{43}{29} \\ -\frac{62}{29} \\ \frac{-7}{29} \end{pmatrix}$ 

Le repère étant orthonormé on peut calculer le produit scalaire :

$\vectt{\text{AX}} \cdot \vectt{\text{BC}} =  \dfrac{172}{29} - \dfrac{186}{29}+ \dfrac{14}{29} = \dfrac{186}{29} - \dfrac{186}{29} = 0$ donc les vecteurs $\vectt{\text{AX}}$ et $\vectt{\text{BC}}$ sont orthogonaux.

D'autre part $\vectt{\text{BX}} = \begin{pmatrix} -\dfrac{72}{29} \\\dfrac{54}{29}\\ -\dfrac{72}{29} \end{pmatrix}$ et on constate que $\vectt{\text{BX}} = \dfrac{18}{29}\vectt{\text{BC}}$.

Conclusion les vecteurs $\vectt{\text{BX}}$ et $\vectt{\text{BC}}$ sont colinéaires ou encore X appartient à la droite (BC).

X est donc le projeté orthogonal de A sur (BC) ; les coordonnées de H sont bien $\left(\dfrac{73}{29}~;~\dfrac{- 4}{29}~;~\dfrac{21}{29} \right)$.
\end{itemize}

Donc le point H est le projeté orthogonal du point A sur la droite (BC).
	\item Le segment [AH] est donc la hauteur issue de A dans le triangle ABC. 
L'aire du triangle ABC est  donc $\mathcal{A}'= \dfrac{\mathrm{BC} \times \mathrm{AH}}{2}$.

$\mathrm{AH}^2 =  \left(-\dfrac{43}{29} \right)^2 + \left(-\dfrac{62}{29}\right)^2 +  \left(-\dfrac{7}{29}\right) ^2 = \dfrac{5742}{29^2} = \dfrac{198}{29} $ 

$\mathrm{BC}^2= (-4)^2+3^2+(-2)^2= 16+9+4=29$

D'où $\mathcal{A}'= \dfrac{ 1}{2} \times \sqrt{\dfrac{198}{29}} \times  \sqrt{29} = \dfrac{\sqrt{198}}{2}=   \dfrac{\sqrt{22 \times 9 }}{2}=\dfrac{3\sqrt{22}}{2}$.

	\item Soit $h$ la distance du point D au plan (ABC). 

$h$ sera la hauteur associé à la base (ABC) dans la pyramide ABCD. 
	
On a donc $V = \dfrac{1}{3} \times \mathcal{A}' \times h$
	
D'où $3 = \dfrac{1}{3} \times  \dfrac{3\sqrt{22}}{2}\times h $ donc $h = \dfrac{3 \times 2}{\sqrt{22}}= \dfrac{3\sqrt{22}}{11}$

La distance du point D au plan (ABC) est donc égale à $\dfrac{3\sqrt{22}}{11}$.
\end{enumerate}

\medskip

\end{document}
